% Options for packages loaded elsewhere
% Options for packages loaded elsewhere
\PassOptionsToPackage{unicode}{hyperref}
\PassOptionsToPackage{hyphens}{url}
\PassOptionsToPackage{dvipsnames,svgnames,x11names}{xcolor}
%
\documentclass[
  letterpaper,
  DIV=11,
  numbers=noendperiod]{scrreprt}
\usepackage{xcolor}
\usepackage{amsmath,amssymb}
\setcounter{secnumdepth}{5}
\usepackage{iftex}
\ifPDFTeX
  \usepackage[T1]{fontenc}
  \usepackage[utf8]{inputenc}
  \usepackage{textcomp} % provide euro and other symbols
\else % if luatex or xetex
  \usepackage{unicode-math} % this also loads fontspec
  \defaultfontfeatures{Scale=MatchLowercase}
  \defaultfontfeatures[\rmfamily]{Ligatures=TeX,Scale=1}
\fi
\usepackage{lmodern}
\ifPDFTeX\else
  % xetex/luatex font selection
\fi
% Use upquote if available, for straight quotes in verbatim environments
\IfFileExists{upquote.sty}{\usepackage{upquote}}{}
\IfFileExists{microtype.sty}{% use microtype if available
  \usepackage[]{microtype}
  \UseMicrotypeSet[protrusion]{basicmath} % disable protrusion for tt fonts
}{}
\makeatletter
\@ifundefined{KOMAClassName}{% if non-KOMA class
  \IfFileExists{parskip.sty}{%
    \usepackage{parskip}
  }{% else
    \setlength{\parindent}{0pt}
    \setlength{\parskip}{6pt plus 2pt minus 1pt}}
}{% if KOMA class
  \KOMAoptions{parskip=half}}
\makeatother
% Make \paragraph and \subparagraph free-standing
\makeatletter
\ifx\paragraph\undefined\else
  \let\oldparagraph\paragraph
  \renewcommand{\paragraph}{
    \@ifstar
      \xxxParagraphStar
      \xxxParagraphNoStar
  }
  \newcommand{\xxxParagraphStar}[1]{\oldparagraph*{#1}\mbox{}}
  \newcommand{\xxxParagraphNoStar}[1]{\oldparagraph{#1}\mbox{}}
\fi
\ifx\subparagraph\undefined\else
  \let\oldsubparagraph\subparagraph
  \renewcommand{\subparagraph}{
    \@ifstar
      \xxxSubParagraphStar
      \xxxSubParagraphNoStar
  }
  \newcommand{\xxxSubParagraphStar}[1]{\oldsubparagraph*{#1}\mbox{}}
  \newcommand{\xxxSubParagraphNoStar}[1]{\oldsubparagraph{#1}\mbox{}}
\fi
\makeatother


\usepackage{longtable,booktabs,array}
\newcounter{none} % for unnumbered tables
\usepackage{calc} % for calculating minipage widths
% Correct order of tables after \paragraph or \subparagraph
\usepackage{etoolbox}
\makeatletter
\patchcmd\longtable{\par}{\if@noskipsec\mbox{}\fi\par}{}{}
\makeatother
% Allow footnotes in longtable head/foot
\IfFileExists{footnotehyper.sty}{\usepackage{footnotehyper}}{\usepackage{footnote}}
\makesavenoteenv{longtable}
\usepackage{graphicx}
\makeatletter
\newsavebox\pandoc@box
\newcommand*\pandocbounded[1]{% scales image to fit in text height/width
  \sbox\pandoc@box{#1}%
  \Gscale@div\@tempa{\textheight}{\dimexpr\ht\pandoc@box+\dp\pandoc@box\relax}%
  \Gscale@div\@tempb{\linewidth}{\wd\pandoc@box}%
  \ifdim\@tempb\p@<\@tempa\p@\let\@tempa\@tempb\fi% select the smaller of both
  \ifdim\@tempa\p@<\p@\scalebox{\@tempa}{\usebox\pandoc@box}%
  \else\usebox{\pandoc@box}%
  \fi%
}
% Set default figure placement to htbp
\def\fps@figure{htbp}
\makeatother


% definitions for citeproc citations
\NewDocumentCommand\citeproctext{}{}
\NewDocumentCommand\citeproc{mm}{%
  \begingroup\def\citeproctext{#2}\cite{#1}\endgroup}
\makeatletter
 % allow citations to break across lines
 \let\@cite@ofmt\@firstofone
 % avoid brackets around text for \cite:
 \def\@biblabel#1{}
 \def\@cite#1#2{{#1\if@tempswa , #2\fi}}
\makeatother
\newlength{\cslhangindent}
\setlength{\cslhangindent}{1.5em}
\newlength{\csllabelwidth}
\setlength{\csllabelwidth}{3em}
\newenvironment{CSLReferences}[2] % #1 hanging-indent, #2 entry-spacing
 {\begin{list}{}{%
  \setlength{\itemindent}{0pt}
  \setlength{\leftmargin}{0pt}
  \setlength{\parsep}{0pt}
  % turn on hanging indent if param 1 is 1
  \ifodd #1
   \setlength{\leftmargin}{\cslhangindent}
   \setlength{\itemindent}{-1\cslhangindent}
  \fi
  % set entry spacing
  \setlength{\itemsep}{#2\baselineskip}}}
 {\end{list}}
\usepackage{calc}
\newcommand{\CSLBlock}[1]{\hfill\break\parbox[t]{\linewidth}{\strut\ignorespaces#1\strut}}
\newcommand{\CSLLeftMargin}[1]{\parbox[t]{\csllabelwidth}{\strut#1\strut}}
\newcommand{\CSLRightInline}[1]{\parbox[t]{\linewidth - \csllabelwidth}{\strut#1\strut}}
\newcommand{\CSLIndent}[1]{\hspace{\cslhangindent}#1}



\setlength{\emergencystretch}{3em} % prevent overfull lines

\providecommand{\tightlist}{%
  \setlength{\itemsep}{0pt}\setlength{\parskip}{0pt}}



 


\KOMAoption{captions}{tableheading}
\makeatletter
\@ifpackageloaded{bookmark}{}{\usepackage{bookmark}}
\makeatother
\makeatletter
\@ifpackageloaded{caption}{}{\usepackage{caption}}
\AtBeginDocument{%
\ifdefined\contentsname
  \renewcommand*\contentsname{Table of contents}
\else
  \newcommand\contentsname{Table of contents}
\fi
\ifdefined\listfigurename
  \renewcommand*\listfigurename{List of Figures}
\else
  \newcommand\listfigurename{List of Figures}
\fi
\ifdefined\listtablename
  \renewcommand*\listtablename{List of Tables}
\else
  \newcommand\listtablename{List of Tables}
\fi
\ifdefined\figurename
  \renewcommand*\figurename{Figure}
\else
  \newcommand\figurename{Figure}
\fi
\ifdefined\tablename
  \renewcommand*\tablename{Table}
\else
  \newcommand\tablename{Table}
\fi
}
\@ifpackageloaded{float}{}{\usepackage{float}}
\floatstyle{ruled}
\@ifundefined{c@chapter}{\newfloat{codelisting}{h}{lop}}{\newfloat{codelisting}{h}{lop}[chapter]}
\floatname{codelisting}{Listing}
\newcommand*\listoflistings{\listof{codelisting}{List of Listings}}
\makeatother
\makeatletter
\makeatother
\makeatletter
\@ifpackageloaded{caption}{}{\usepackage{caption}}
\@ifpackageloaded{subcaption}{}{\usepackage{subcaption}}
\makeatother
\usepackage{bookmark}
\IfFileExists{xurl.sty}{\usepackage{xurl}}{} % add URL line breaks if available
\urlstyle{same}
\hypersetup{
  pdftitle={Rekayasa Multidisiplin di Era Kecerdasan Sisttem},
  pdfauthor={Armein Z. R. Langi},
  colorlinks=true,
  linkcolor={blue},
  filecolor={Maroon},
  citecolor={Blue},
  urlcolor={Blue},
  pdfcreator={LaTeX via pandoc}}


\title{Rekayasa Multidisiplin di Era Kecerdasan Sisttem}
\usepackage{etoolbox}
\makeatletter
\providecommand{\subtitle}[1]{% add subtitle to \maketitle
  \apptocmd{\@title}{\par {\large #1 \par}}{}{}
}
\makeatother
\subtitle{Triune-Intelligence Smart Engineering Level 5}
\author{Armein Z. R. Langi}
\date{2026-04-20}
\begin{document}
\maketitle

\renewcommand*\contentsname{Table of contents}
{
\setcounter{tocdepth}{2}
\tableofcontents
}

\bookmarksetup{startatroot}

\chapter*{Prakata}\label{prakata}
\addcontentsline{toc}{chapter}{Prakata}

\markboth{Prakata}{Prakata}

Monograf ini disusun untuk merumuskan secara sistematis perkembangan
konseptual dalam paradigma \textbf{Triune-Intelligence Smart Engineering
(TISE)} hingga mencapai \textbf{TISE level 5}, yaitu level di mana TISE
hadir sebagai pusat keunggulan rekayasa bagi penyelesaian persoalan
kemanusiaan. Jika rekayasa klasik berhasil besar dalam membangun
infrastruktur peradaban material, maka masyarakat abad ke-21 menghadapi
tuntutan baru: bagaimana membantu manusia hidup secara aktif, produktif,
dan bermakna melalui konfigurasi sosial yang saling menopang misi.

Monograf ini berangkat dari keyakinan bahwa rekayasa perlu bergerak
melampaui batas-batas sektoralnya. Rekayasa tidak lagi hanya berbicara
tentang mesin fisik, proses industri, atau sistem komputasi, tetapi juga
tentang bagaimana manusia, institusi, teknologi, narasi, dan nilai dapat
diorganisasikan menjadi solusi yang hidup. Karena itu, monograf ini
menelusuri evolusi abstraksi rekayasa dari \textbf{core engine} pada
TISE level 0, menuju \textbf{smart engine/mesin PSKVE} pada TISE level
1, \textbf{mesin agentik PUDAL} pada TISE level 2, \textbf{teater
identitas naratif} pada TISE level 3, \textbf{DNA desain meta-TISE} pada
TISE level 4, hingga \textbf{Matrix of Mission} sebagai mesin
konfigurasi stakeholder pada TISE level 5.

Secara khusus, monograf ini menempatkan \textbf{MOS-7} sebagai siklus
operasional TISE level 5 dan \textbf{Matrix of Mission} sebagai mesin
utama untuk membentuk serta mengevaluasi konfigurasi stakeholder
bermisi. Melalui kedua perangkat ini, TISE level 5 dirumuskan sebagai
bentuk kelembagaan rekayasa yang mampu memobilisasi potensi manusia yang
berlimpah untuk memecahkan persoalan kemanusiaan dengan prinsip
\textbf{dari manusia untuk manusia}.

Sebagai monograf, tulisan ini tidak dimaksudkan sebagai laporan proyek
atau manual teknis semata, melainkan sebagai fondasi konseptual yang
dapat menopang riset, pendidikan, desain artefak, pembentukan pusat
keunggulan, dan pengembangan metodologi lanjutan. Lebih jauh lagi,
monograf ini diposisikan sebagai sebuah \textbf{Buku DNA}, yaitu dokumen
acuan yang dapat dibaca, ditafsirkan, dan dijalankan oleh para
rekayasawan, stakeholder, dan agents dalam bagian dan peran mereka
masing-masing. Dengan demikian, monograf ini diharapkan menjadi titik
tolak bagi pembentukan bidang kajian rekayasa yang lebih manusiawi,
lebih reflektif, dan lebih relevan bagi tantangan abad ke-21, sekaligus
menjadi naskah pengarah bagi perwujudan nyata TISE level 5.

\bookmarksetup{startatroot}

\chapter*{Abstrak}\label{abstrak}
\addcontentsline{toc}{chapter}{Abstrak}

\markboth{Abstrak}{Abstrak}

Monograf ini merumuskan \textbf{TISE level 5} sebagai level perkembangan
rekayasa di mana paradigma \textbf{Triune-Intelligence Smart Engineering
(TISE)} hadir dalam bentuk \textbf{pusat keunggulan rekayasa} untuk
memecahkan persoalan kemanusiaan. Argumen utamanya adalah bahwa
persoalan manusia abad ke-21 tidak lagi memadai jika dipahami hanya
sebagai persoalan energi, material, dan informasi, melainkan perlu
dipandang sebagai persoalan konfigurasi stakeholder bermisi, distribusi
benefit, legitimasi regulatif, operasi solusi, dan mobilisasi potensi
manusia.

Untuk menjelaskan pergeseran tersebut, monograf ini menyusun hirarki
level rekayasa: \textbf{TISE level 0} sebagai rekayasa tradisional
berbasis core engine; \textbf{TISE level 1} sebagai smart engineering
berbasis smart engine/mesin PSKVE; \textbf{TISE level 2} sebagai
rekayasa agentik berbasis mesin PUDAL; \textbf{TISE level 3} sebagai
rekayasa pementasan identitas naratif; \textbf{TISE level 4} sebagai
rekayasa meta-artefak berbasis DNA desain; dan \textbf{TISE level 5}
sebagai rekayasa konfigurasi stakeholder berbasis Matrix of Mission.

Dalam monograf ini, \textbf{Matrix of Mission} dirumuskan sebagai mesin
utama TISE level 5, sedangkan \textbf{MOS-7} dirumuskan sebagai siklus
operasional untuk mendengar keluhan manusia, memodelkan stakeholder,
mendesain konfigurasi solusi, membangun komitmen, membangun artefak,
mengoperasionalkan solusi, dan mengevaluasi integritas serta dampaknya.
Dari sudut pandang literatur, posisi ini dapat dipahami sebagai
perluasan atas tradisi PSS, stakeholder motivation matrix, dan
multi-stakeholder requirements (Morelli 2002; Vezzoli et al. 2014;
Berkovich et al. 2014; Sholihah et al. 2019). Dengan kerangka tersebut,
TISE level 5 dipahami bukan sekadar sebagai metode desain, melainkan
sebagai bentuk kelembagaan yang memobilisasi manusia untuk menolong
manusia lain, sekaligus membuka lapangan kerja baru dalam pelayanan
kemanusiaan, operasi solusi, pendampingan, dan fasilitasi komunitas.

Kontribusi utama monograf ini meliputi: (1) perumusan ontologi
bertingkat bagi evolusi rekayasa dari core engine ke Matrix of Mission;
(2) generalisasi konsep energi menjadi \textbf{energon}; (3) definisi
formal atas PSKVE, PUDAL, stakeholder bermisi, Matrix of Mission, dan
TISE level 5; serta (4) proposisi konseptual dan agenda riset bagi
pengembangan TISE level 5 sebagai bidang ilmu dan praktik. Dengan
demikian, monograf ini menawarkan dasar konseptual bagi rekayasa yang
lebih integratif, partisipatif, agentik, dan manusiawi.

\textbf{Kata kunci:} TISE level 5, monograf rekayasa, Matrix of Mission,
MOS-7, stakeholder bermisi, energon, PSKVE, PUDAL, pusat keunggulan,
rekayasa kemanusiaan.

\bookmarksetup{startatroot}

\chapter*{Bagian I. Pengantar Monograf dan Posisi TISE level
5}\label{bagian-i.-pengantar-monograf-dan-posisi-tise-level-5}
\addcontentsline{toc}{chapter}{Bagian I. Pengantar Monograf dan Posisi
TISE level 5}

\markboth{Bagian I. Pengantar Monograf dan Posisi TISE level 5}{Bagian
I. Pengantar Monograf dan Posisi TISE level 5}

\bookmarksetup{startatroot}

\chapter*{Bab 1. Pendahuluan}\label{bab-1.-pendahuluan}
\addcontentsline{toc}{chapter}{Bab 1. Pendahuluan}

\markboth{Bab 1. Pendahuluan}{Bab 1. Pendahuluan}

\section*{Latar Belakang}\label{latar-belakang}
\addcontentsline{toc}{section}{Latar Belakang}

\markright{Latar Belakang}

Sepanjang sejarah peradaban modern, rekayasa telah terbukti menjadi
salah satu kekuatan utama yang memungkinkan manusia mengatasi berbagai
persoalan hidupnya. Rekayasa telah melahirkan pemukiman yang lebih
layak, gedung dan tempat kerja yang menopang produktivitas, jalan raya
yang menghubungkan wilayah, transportasi udara yang mempercepat
mobilitas, sistem pangan yang mendukung populasi besar, infrastruktur
energi yang memungkinkan aktivitas ekonomi, teknologi komunikasi yang
menghubungkan manusia lintas jarak, serta industri yang menghasilkan
barang dan layanan dalam skala besar. Dalam pengertian ini, rekayasa
telah menjadi salah satu fondasi terpenting bagi keberhasilan peradaban
dalam menanggapi kebutuhan-kebutuhan dasar manusia.

Namun, pada abad ke-21, persoalan manusia tidak lagi berhenti pada
penyediaan kebutuhan dasar material dan infrastruktur. Persoalan mulai
bergeser ke arah kebutuhan akan kehidupan yang \textbf{aktif, produktif,
dan bermakna}. Manusia tidak cukup hanya memiliki rumah, jalan, listrik,
makanan, atau alat komunikasi. Manusia ingin hidup sebagai pribadi yang
memiliki arah, mampu berkarya, dapat berkontribusi, dan merasa bahwa
hidupnya berarti. Kebutuhan ini semakin nyata dalam masyarakat yang
kompleks, terdigitalisasi, dan saling bergantung, di mana keberhasilan
seseorang sangat dipengaruhi oleh kemampuan untuk memahami misi dirinya
sendiri, memahami misi orang lain, dan membangun konfigurasi sosial yang
saling menopang.

Dengan demikian, kebutuhan manusia modern semakin bersifat
\textbf{misioner}. Yang dicari bukan sekadar survival, melainkan
kemampuan menjalankan misi kehidupan secara aktif dan produktif bersama
orang lain. Dalam konteks ini, persoalan kemanusiaan sering kali muncul
bukan hanya karena kurangnya teknologi, tetapi karena tidak terbangunnya
konfigurasi stakeholder yang saling memahami, saling menopang, dan
bersama-sama membentuk solusi yang hidup. Karena itu, dibutuhkan
lompatan paradigma: bukan hanya menggunakan engineering untuk
menghasilkan artefak, tetapi juga menjadikan engineering sebagai cara
masyarakat membentuk dirinya sendiri.

Di sinilah muncul kebutuhan akan dua gerakan yang saling terkait, yaitu
\textbf{memasyarakatkan engineering} dan \textbf{mengengineeringkan
masyarakat}. Memasyarakatkan engineering berarti menjadikan cara
berpikir rekayasa, cara memodelkan masalah, cara membangun solusi, dan
cara mengevaluasi sistem sebagai bagian dari budaya masyarakat luas.
Sebaliknya, mengengineeringkan masyarakat berarti memperlakukan
masyarakat, relasi sosial, organisasi, dan konfigurasi stakeholder
sebagai sesuatu yang dapat direkayasa secara sadar agar menjadi lebih
efektif, adil, produktif, dan bermakna. Dengan kata lain, masyarakat
perlu belajar menjadi engineers dalam arti luas, sementara engineers
perlu semakin sungguh-sungguh mengarahkan ilmunya untuk mengatasi
persoalan masyarakat.

Transformasi semacam ini tidak dapat dijalankan hanya dengan niat baik
atau teknologi semata. Diperlukan metodologi yang efektif untuk
mendengar persoalan manusia, memodelkan stakeholder, membangun komitmen,
merekayasa artefak, mengoperasionalkan solusi, dan mengevaluasi
dampaknya secara berkelanjutan. Kebutuhan inilah yang melatarbelakangi
pengembangan paradigma \textbf{Triune-Intelligence Smart Engineering
(TISE)}.

\section*{Rumusan Masalah}\label{rumusan-masalah}
\addcontentsline{toc}{section}{Rumusan Masalah}

\markright{Rumusan Masalah}

Monograf ini berangkat dari beberapa pertanyaan utama:

\begin{enumerate}
\def\labelenumi{\arabic{enumi}.}
\tightlist
\item
  Bagaimana TISE level 5 dapat dirumuskan sebagai pusat keunggulan
  rekayasa?
\item
  Bagaimana masalah kemanusiaan dipahami sebagai objek rekayasa dalam
  paradigma TISE?
\item
  Bagaimana MOS-7 dapat diadaptasi menjadi siklus operasional TISE level
  5?
\item
  Bagaimana stakeholder diberdayakan untuk menjalankan misi
  masing-masing dalam pembangunan artefak solusi?
\item
  Bagaimana integritas, etika, dan keberlanjutan solusi dijaga?
\end{enumerate}

\section*{Tujuan Monograf}\label{tujuan-monograf}
\addcontentsline{toc}{section}{Tujuan Monograf}

\markright{Tujuan Monograf}

Monograf ini bertujuan:

\begin{enumerate}
\def\labelenumi{\arabic{enumi}.}
\tightlist
\item
  merumuskan definisi dan posisi konseptual TISE level 5 dalam evolusi
  rekayasa;
\item
  menjelaskan hubungan antara TISE level 0, TISE level 1, TISE level 2,
  TISE level 3, TISE level 4, dan TISE level 5;
\item
  merumuskan Matrix of Mission sebagai mesin utama konfigurasi
  stakeholder;
\item
  merumuskan MOS-7 sebagai \emph{mission operating cycle} bagi pusat
  keunggulan rekayasa kemanusiaan;
\item
  menawarkan ontologi, definisi formal, dan proposisi ilmiah bagi
  pengembangan TISE level 5;
\item
  memberi dasar bagi penelitian, pengajaran, pembentukan pusat
  keunggulan, dan implementasi kelembagaan TISE level 5.
\end{enumerate}

\section*{Kontribusi Monograf}\label{kontribusi-monograf}
\addcontentsline{toc}{section}{Kontribusi Monograf}

\markright{Kontribusi Monograf}

Kontribusi utama monograf ini adalah menggeser TISE dari sekadar
paradigma atau metodologi menjadi \textbf{institusi rekayasa yang
operasional}, yaitu pusat keunggulan yang mampu memobilisasi banyak agen
TISE dan banyak stakeholder di sekitar misi-misi kemanusiaan.

Secara lebih rinci, kontribusi monograf ini dapat diringkas sebagai
berikut:

\begin{enumerate}
\def\labelenumi{\arabic{enumi}.}
\tightlist
\item
  menyusun \textbf{hirarki evolusi rekayasa} dari TISE level 0 hingga
  TISE level 5;
\item
  memperkenalkan konsep \textbf{energon} sebagai generalisasi energi
  bagi TISE level 1;
\item
  merumuskan \textbf{smart engine/mesin PSKVE} sebagai abstraksi utama
  TISE level 1;
\item
  menempatkan \textbf{PUDAL}, \textbf{identitas naratif}, dan
  \textbf{DNA desain} sebagai level-level abstraksi lanjutan dalam
  perkembangan TISE;
\item
  merumuskan \textbf{Matrix of Mission} sebagai mesin utama konfigurasi
  stakeholder pada TISE level 5;
\item
  menempatkan \textbf{MOS-7} sebagai siklus operasional bagi pusat
  keunggulan TISE level 5;
\item
  merumuskan agenda riset untuk pengembangan TISE level 5 sebagai bidang
  ilmu dan praktik.
\end{enumerate}

\section*{Posisi Monograf dalam
Literatur}\label{posisi-monograf-dalam-literatur}
\addcontentsline{toc}{section}{Posisi Monograf dalam Literatur}

\markright{Posisi Monograf dalam Literatur}

Monograf ini berada pada persimpangan beberapa rumpun literatur.
Pertama, literatur \textbf{Product-Service Systems (PSS)} menegaskan
bahwa solusi kontemporer semakin hadir sebagai kombinasi produk dan
layanan yang terintegrasi untuk memenuhi kebutuhan pengguna, bukan
sebagai produk tunggal. Kedua, literatur \textbf{service design} dan
\textbf{system design tools} telah memperkenalkan alat-alat visual untuk
memetakan aktor, relasi, dan motivasi dalam sistem layanan. Ketiga,
literatur \textbf{multi-stakeholder requirements} menunjukkan bahwa
solusi yang bernilai sering bergantung pada kemampuan menangkap dan
mengoordinasikan kebutuhan banyak pihak sekaligus. Keempat, literatur
\textbf{stakeholder engagement and management} menekankan bahwa kualitas
solusi sangat dipengaruhi oleh cara para aktor didefinisikan,
dilibatkan, dan dihubungkan dalam keseluruhan siklus sistem.

Dalam konteks itu, monograf ini menempatkan \textbf{Matrix of Mission}
sebagai pengembangan dari tradisi \textbf{stakeholder motivation matrix}
dalam literatur PSS dan service design (Morelli 2002; Vezzoli et al.
2014; Giordano et al. 2018). Namun, pengembangan yang diusulkan di sini
bergerak melampaui motivasi menuju \textbf{misi}, \textbf{benefit},
\textbf{kontribusi silang}, \textbf{tipologi stakeholder} (USER, SOURCE,
REGULATOR, PROVIDER), dan \textbf{operasi siklik} melalui MOS-7. Dengan
demikian, monograf ini tidak hanya mengusulkan alat visual baru, tetapi
juga sebuah kerangka ontologis dan kelembagaan baru bagi rekayasa
kemanusiaan.

\bookmarksetup{startatroot}

\chapter*{Bab 2. Tinjauan Pustaka dan Dasar
Konseptual}\label{bab-2.-tinjauan-pustaka-dan-dasar-konseptual}
\addcontentsline{toc}{chapter}{Bab 2. Tinjauan Pustaka dan Dasar
Konseptual}

\markboth{Bab 2. Tinjauan Pustaka dan Dasar Konseptual}{Bab 2. Tinjauan
Pustaka dan Dasar Konseptual}

\section*{Product-Service Systems sebagai Latar
Penting}\label{product-service-systems-sebagai-latar-penting}
\addcontentsline{toc}{section}{Product-Service Systems sebagai Latar
Penting}

\markright{Product-Service Systems sebagai Latar Penting}

Salah satu landasan penting bagi monograf ini adalah literatur
\textbf{Product-Service Systems (PSS)}. Nicola Morelli menekankan
perlunya pergeseran dari produksi barang menuju penyediaan solusi
sistemik yang terdiri atas kombinasi produk dan layanan (Morelli 2002).
Dalam pengertian ini, desain tidak lagi hanya menangani bentuk artefak,
tetapi juga hubungan antaraktor, proses penggunaan, dan struktur
penyampaian nilai. Literatur PSS selanjutnya juga menegaskan bahwa PSS
merupakan bundel elemen fisik dan layanan yang diintegrasikan untuk
menyelesaikan masalah pengguna (Vezzoli et al. 2014; Berkovich et al.
2014).

\section*{Stakeholder Motivation Matrix dan Akar Matrix of
Mission}\label{stakeholder-motivation-matrix-dan-akar-matrix-of-mission}
\addcontentsline{toc}{section}{Stakeholder Motivation Matrix dan Akar
Matrix of Mission}

\markright{Stakeholder Motivation Matrix dan Akar Matrix of Mission}

Literatur desain untuk PSS telah mengembangkan alat-alat representasi
sistem seperti \emph{system map}, \emph{interaction storyboard}, dan
\textbf{stakeholder motivation matrix} (Vezzoli et al. 2014). Dalam
literatur tersebut, \emph{motivation matrix} dijelaskan sebagai tabel
dua-arah yang memungkinkan perancang melihat motivasi tiap aktor,
kontribusi yang diberikan kepada kemitraan, \emph{expected benefits},
serta potensi sinergi atau konflik dengan aktor lain (Vezzoli et al.
2014).

Monograf ini mengakui secara eksplisit bahwa \textbf{Matrix of Mission}
berada dalam garis pengembangan konseptual dari tradisi tersebut
(Morelli 2002; Vezzoli et al. 2014). Namun, Matrix of Mission yang
diusulkan di sini memperluas fokus \emph{motivation matrix}. Jika
\emph{motivation matrix} terutama bertanya mengapa aktor mau terlibat
dan apa yang mereka kontribusikan atau harapkan, maka Matrix of Mission
menempatkan \textbf{misi} sebagai pusat diagonal utama dan
menghubungkannya dengan benefit, kontribusi antarpihak, keberlanjutan
operasi, serta evaluasi kelayakan konfigurasi stakeholder. Karena itu,
Matrix of Mission bukan hanya alat motivasional, melainkan mesin
konfigurasi sosial-operasional.

\section*{Multi-Stakeholder Requirements dan Konfigurasi
Solusi}\label{multi-stakeholder-requirements-dan-konfigurasi-solusi}
\addcontentsline{toc}{section}{Multi-Stakeholder Requirements dan
Konfigurasi Solusi}

\markright{Multi-Stakeholder Requirements dan Konfigurasi Solusi}

Literatur PSS yang lebih baru juga menunjukkan bahwa sistem nyata
biasanya dibentuk oleh \textbf{konsorsium multi-stakeholder}, dan
keberhasilan desain sangat bergantung pada kemampuan menangkap serta
mewujudkan kebutuhan banyak pihak secara serempak (Sholihah et al.
2019). Dengan demikian, fokus pada kebutuhan aktor tunggal tidak lagi
memadai. Solusi harus dipahami sebagai hasil negosiasi, integrasi, dan
pengaturan ketergantungan antarpihak.

\section*{Stakeholder Engagement, Stakeholder Map, dan Rekayasa
Kemanusiaan}\label{stakeholder-engagement-stakeholder-map-dan-rekayasa-kemanusiaan}
\addcontentsline{toc}{section}{Stakeholder Engagement, Stakeholder Map,
dan Rekayasa Kemanusiaan}

\markright{Stakeholder Engagement, Stakeholder Map, dan Rekayasa
Kemanusiaan}

Literatur service design menunjukkan bahwa \textbf{stakeholder map}
berperan sebagai alat percakapan untuk membangun layanan yang lebih
\emph{people-led} dan \emph{citizen-centred} (Giordano et al. 2018).
Literatur lain tentang stakeholder engagement dalam inovasi juga
menekankan bahwa ekosistem inovasi modern terdiri dari konstelasi
stakeholder yang heterogen, dengan tujuan, motif, dan kapabilitas yang
berbeda (Lievens and Blazevic 2021).

\section*{Posisi Kebaruan Monograf}\label{posisi-kebaruan-monograf}
\addcontentsline{toc}{section}{Posisi Kebaruan Monograf}

\markright{Posisi Kebaruan Monograf}

Kebaruan utama monograf ini terletak pada lima hal. Pertama, monograf
ini menyusun \textbf{hirarki evolusi rekayasa} dari TISE level 0 hingga
TISE level 5. Kedua, monograf ini memperkenalkan \textbf{energon}
sebagai generalisasi energi. Ketiga, monograf ini menempatkan
\textbf{PSKVE} dan \textbf{PUDAL} ke dalam kerangka rekayasa bertingkat.
Keempat, monograf ini mengembangkan \textbf{Matrix of Mission} dari
tradisi stakeholder motivation matrix menjadi mesin konfigurasi
stakeholder bermisi. Kelima, monograf ini menghubungkan semuanya dengan
\textbf{MOS-7} sebagai siklus operasional bagi pusat keunggulan rekayasa
kemanusiaan.

\bookmarksetup{startatroot}

\chapter*{Bagian II. Evolusi Rekayasa Menuju TISE level
5}\label{bagian-ii.-evolusi-rekayasa-menuju-tise-level-5}
\addcontentsline{toc}{chapter}{Bagian II. Evolusi Rekayasa Menuju TISE
level 5}

\markboth{Bagian II. Evolusi Rekayasa Menuju TISE level 5}{Bagian II.
Evolusi Rekayasa Menuju TISE level 5}

\bookmarksetup{startatroot}

\chapter*{Bab 3. Evolusi Paradigma Rekayasa: dari TISE level 0 hingga
TISE level
5}\label{bab-3.-evolusi-paradigma-rekayasa-dari-tise-level-0-hingga-tise-level-5}
\addcontentsline{toc}{chapter}{Bab 3. Evolusi Paradigma Rekayasa: dari
TISE level 0 hingga TISE level 5}

\markboth{Bab 3. Evolusi Paradigma Rekayasa: dari TISE level 0 hingga
TISE level 5}{Bab 3. Evolusi Paradigma Rekayasa: dari TISE level 0
hingga TISE level 5}

\section*{TISE level 0: Rekayasa Tradisional Berbasis Core
Engine}\label{tise-level-0-rekayasa-tradisional-berbasis-core-engine}
\addcontentsline{toc}{section}{TISE level 0: Rekayasa Tradisional
Berbasis Core Engine}

\markright{TISE level 0: Rekayasa Tradisional Berbasis Core Engine}

Pada \textbf{TISE level 0}, rekayasa dapat dipahami sebagai kegiatan
\textbf{mengelola energi menjadi kerja} melalui artefak-artefak yang
dapat diabstraksikan sebagai \textbf{mesin inti} (\emph{core engine}).
Mesin, dalam arti luas, adalah konfigurasi instrumen yang mengubah,
menyalurkan, menyimpan, mengendalikan, atau menstabilkan energi agar
suatu kerja tertentu dapat dicapai. Di balik keragaman disiplin
rekayasa, terdapat pola umum: rekayasa adalah desain dan implementasi
mesin energi untuk memecahkan masalah yang terkait dengan \textbf{beban
energi, material, dan informasi (EMI)}.

Pada level ini, mesin dikendalikan terutama melalui mekanisme sederhana
seperti \textbf{on-off} atau melalui satu besaran kendali dominan yang
bersifat satu dimensi, misalnya percepatan atau perlambatan. Secara
fungsional terdapat dua jenis mesin: \textbf{mesin statik} yang
dirancang untuk menjaga kestabilan struktur atau keadaan diam, dan
\textbf{mesin dinamik} yang dirancang untuk mengatur perpindahan,
kecepatan, percepatan, aliran, atau perubahan keadaan.

\section*{TISE level 1: Smart Engineering Berbasis Smart Engine atau
Mesin
PSKVE}\label{tise-level-1-smart-engineering-berbasis-smart-engine-atau-mesin-pskve}
\addcontentsline{toc}{section}{TISE level 1: Smart Engineering Berbasis
Smart Engine atau Mesin PSKVE}

\markright{TISE level 1: Smart Engineering Berbasis Smart Engine atau
Mesin PSKVE}

Pada level ini, logika kendali diperluas dari satu besaran dominan ke
\textbf{lima dimensi}, yaitu \textbf{P, S, K, V, dan E} (\emph{Product,
Service, Knowledge, Value, Environment}). Karena itu, sistem kendalinya
memerlukan \textbf{lima derajat kebebasan}, sehingga dibutuhkan
komputer, algoritma pengendalian, dan kemampuan komputasional yang lebih
tinggi.

TISE level 1 adalah perluasan dari TISE level 0. Jika pada TISE level 0
fokus utama berada pada konversi energi, material, dan informasi, maka
pada TISE level 1 fokusnya diperluas pada konversi berbagai jenis
energon yang hadir dalam kehidupan manusia, organisasi, dan masyarakat.
TISE level 1 karena itu dapat didefinisikan sebagai \textbf{rekayasa
multidisiplin untuk membangun instrumen konversi energon dalam
konfigurasi mesin PSKVE}.

\section*{TISE level 2: Rekayasa Agentik Berbasis Mesin
PUDAL}\label{tise-level-2-rekayasa-agentik-berbasis-mesin-pudal}
\addcontentsline{toc}{section}{TISE level 2: Rekayasa Agentik Berbasis
Mesin PUDAL}

\markright{TISE level 2: Rekayasa Agentik Berbasis Mesin PUDAL}

Pada level ini, mesin PSKVE tidak lagi cukup dikendalikan hanya oleh
komputer dan algoritma biasa. Ia dikendalikan oleh \textbf{mesin cerdas
PUDAL} (\emph{Perception, Understanding, Decision, Action, Learning}).
Di sini \textbf{Triune Intelligence} berinteraksi untuk menjalankan
fungsi kendali statik maupun dinamik. Prompts berfungsi sebagai
mekanisme komunikasi antaragents.

\section*{TISE level 3: Rekayasa Pementasan Identitas
Naratif}\label{tise-level-3-rekayasa-pementasan-identitas-naratif}
\addcontentsline{toc}{section}{TISE level 3: Rekayasa Pementasan
Identitas Naratif}

\markright{TISE level 3: Rekayasa Pementasan Identitas Naratif}

Pada level ini, pementasan \textbf{identitas naratif} menjadi penggerak
utama atau \textbf{energon utama} dari mesin-mesin inti, PSKVE, dan
PUDAL. Asumsi dasarnya adalah bahwa setiap orang memiliki misi di dunia
ini, yaitu mendeliver identitas naratifnya ke dunia sebagai tugas
hidupnya. Artefak rekayasa pada level ini dipahami sebagai
\textbf{teater megah} tempat identitas naratif dipentaskan sehingga
memberi inspirasi atau solusi yang dibutuhkan audiens.

\section*{TISE level 4: Rekayasa Meta-TISE Berbasis DNA
Desain}\label{tise-level-4-rekayasa-meta-tise-berbasis-dna-desain}
\addcontentsline{toc}{section}{TISE level 4: Rekayasa Meta-TISE Berbasis
DNA Desain}

\markright{TISE level 4: Rekayasa Meta-TISE Berbasis DNA Desain}

Pada level ini, rekayasa bergerak ke ranah \textbf{meta-TISE}, yaitu
rekayasa terhadap artefak TISE itu sendiri agar menjadi platform untuk
melahirkan replika atau ``anak artefak''. Pada level ini digunakan
metafora biologis berupa \textbf{DNA desain}. Sebuah sistem terdiri dari
berbagai agents yang menggunakan \textbf{buku skills} yang sama, tetapi
masing-masing berfokus pada bagian tertentu dari keseluruhan sistem.

\section*{TISE level 5: Rekayasa Konfigurasi Stakeholder Berbasis Matrix
of
Mission}\label{tise-level-5-rekayasa-konfigurasi-stakeholder-berbasis-matrix-of-mission}
\addcontentsline{toc}{section}{TISE level 5: Rekayasa Konfigurasi
Stakeholder Berbasis Matrix of Mission}

\markright{TISE level 5: Rekayasa Konfigurasi Stakeholder Berbasis
Matrix of Mission}

Pada level ini, rekayasa mencapai tingkat kelembagaan dan sosial yang
lebih matang. Fokus utamanya bukan lagi mesin inti, mesin PSKVE, PUDAL,
identitas naratif, atau DNA desain, tetapi \textbf{organisasi,
konsorsium, atau komunitas stakeholder} yang dibentuk untuk memecahkan
persoalan kemanusiaan. Pada level ini, artefak rekayasa utama adalah
\textbf{konfigurasi stakeholder} itu sendiri dan mesinnya adalah
\textbf{Matrix of Mission}.

\section*{Tabel Perbandingan TISE level
0--5}\label{tabel-perbandingan-tise-level-05}
\addcontentsline{toc}{section}{Tabel Perbandingan TISE level 0--5}

\markright{Tabel Perbandingan TISE level 0--5}

Untuk memperjelas argumen tentang evolusi abstraksi rekayasa,
\textbf{Tabel 3.1} merangkum perkembangan level rekayasa dari TISE level
0 hingga TISE level 5. Tabel ini perlu dibaca sebagai peta konseptual
yang menunjukkan pergeseran fokus, jenis mesin, mekanisme kendali,
artefak utama, dan tujuan rekayasa pada setiap level.

{\def\LTcaptype{none} % do not increment counter
\begin{longtable}[]{@{}
  >{\raggedright\arraybackslash}p{(\linewidth - 14\tabcolsep) * \real{0.1250}}
  >{\raggedright\arraybackslash}p{(\linewidth - 14\tabcolsep) * \real{0.1250}}
  >{\raggedright\arraybackslash}p{(\linewidth - 14\tabcolsep) * \real{0.1250}}
  >{\raggedright\arraybackslash}p{(\linewidth - 14\tabcolsep) * \real{0.1250}}
  >{\raggedright\arraybackslash}p{(\linewidth - 14\tabcolsep) * \real{0.1250}}
  >{\raggedright\arraybackslash}p{(\linewidth - 14\tabcolsep) * \real{0.1250}}
  >{\raggedright\arraybackslash}p{(\linewidth - 14\tabcolsep) * \real{0.1250}}
  >{\raggedright\arraybackslash}p{(\linewidth - 14\tabcolsep) * \real{0.1250}}@{}}
\toprule\noalign{}
\begin{minipage}[b]{\linewidth}\raggedright
Level
\end{minipage} & \begin{minipage}[b]{\linewidth}\raggedright
Nama Level
\end{minipage} & \begin{minipage}[b]{\linewidth}\raggedright
Fokus Rekayasa
\end{minipage} & \begin{minipage}[b]{\linewidth}\raggedright
Jenis Mesin
\end{minipage} & \begin{minipage}[b]{\linewidth}\raggedright
Energon / Sumber Kerja Utama
\end{minipage} & \begin{minipage}[b]{\linewidth}\raggedright
Mekanisme Kendali
\end{minipage} & \begin{minipage}[b]{\linewidth}\raggedright
Artefak Utama
\end{minipage} & \begin{minipage}[b]{\linewidth}\raggedright
Tujuan Utama
\end{minipage} \\
\midrule\noalign{}
\endhead
\bottomrule\noalign{}
\endlastfoot
0 & Rekayasa tradisional & Konversi EMI menjadi kerja atau bentuk
tersimpan/tertransmisikan & Core Engine & Energi fisik, material,
informasi & On-off, besaran satu dimensi, kendali statik/dinamik klasik
& Mesin, struktur, proses, perangkat sektoral & Menyelesaikan persoalan
sektoral secara stabil, efisien, dan andal \\
1 & Smart Engineering & Konversi energon menjadi PSKVE & Smart Engine /
Mesin PSKVE & Energon 1--5 & Kendali multidimensi 5 derajat kebebasan &
Lingkungan kerja cerdas: stasiun, jalan, kendaraan & Menghasilkan
produk, layanan, pengetahuan, nilai, dan lingkungan yang memecahkan
masalah \\
2 & Rekayasa agentik & Pengendalian agentik atas mesin PSKVE & Mesin
PUDAL & Triune Intelligence yang bekerja melalui agen & PUDAL dan
prompts & Sistem agentik yang mengorkestrasi PSKVE & Menutup gap secara
adaptif dan cerdas \\
3 & Rekayasa naratif & Pementasan identitas naratif sebagai energon
utama & Teater Naratif & Identitas naratif, misi hidup, agensi, makna &
Interaksi naratif dan prompting reflektif & Teater kehidupan /
lingkungan naratif & Memberdayakan stakeholder untuk mendeliver
misinya \\
4 & Rekayasa meta-TISE & Rekayasa meta-artefak yang dapat melahirkan
artefak baru & Mesin DNA Desain / Meta-TISE & Buku skills dan pola
desain & Konsistensi DNA desain & Platform TISE yang dapat menghasilkan
replika atau anak artefak & Membangun platform rekayasa yang
reproduktif \\
5 & Rekayasa konfigurasi stakeholder & Rekayasa konfigurasi stakeholder
bermisi & Matrix of Mission & Potensi manusia, misi stakeholder,
komitmen kolektif & Matrix of Mission dan MOS-7 & Organisasi,
konsorsium, komunitas stakeholder, pusat keunggulan & Memobilisasi
manusia untuk memecahkan persoalan kemanusiaan \\
\end{longtable}
}

\section*{Definisi Ringkas TISE level
5}\label{definisi-ringkas-tise-level-5}
\addcontentsline{toc}{section}{Definisi Ringkas TISE level 5}

\markright{Definisi Ringkas TISE level 5}

\textbf{TISE level 5} adalah level rekayasa di mana TISE hadir sebagai
pusat keunggulan yang memobilisasi stakeholder bermisi, agen-agen TISE,
dan artefak cerdas untuk memecahkan persoalan kemanusiaan melalui
konfigurasi stakeholder, pengoperasian MOS-7, dan penggunaan Matrix of
Mission secara berulang, adaptif, dan bermakna.

\bookmarksetup{startatroot}

\chapter*{Bab 4. Ontologi dan Definisi
Formal}\label{bab-4.-ontologi-dan-definisi-formal}
\addcontentsline{toc}{chapter}{Bab 4. Ontologi dan Definisi Formal}

\markboth{Bab 4. Ontologi dan Definisi Formal}{Bab 4. Ontologi dan
Definisi Formal}

\section*{Kebutuhan akan Definisi
Formal}\label{kebutuhan-akan-definisi-formal}
\addcontentsline{toc}{section}{Kebutuhan akan Definisi Formal}

\markright{Kebutuhan akan Definisi Formal}

Agar TISE level 5 dapat berkembang sebagai paradigma ilmiah, ia tidak
cukup dijelaskan secara metaforis atau naratif saja. Diperlukan ontologi
dan definisi formal yang memungkinkan istilah-istilah utamanya dipakai
secara konsisten dalam penelitian, pengajaran, desain sistem, dan
evaluasi implementasi.

\section*{Definisi Energon}\label{definisi-energon}
\addcontentsline{toc}{section}{Definisi Energon}

\markright{Definisi Energon}

\textbf{Energon} adalah setiap entitas, kondisi, status, kapasitas, atau
sumber daya yang memiliki kemampuan untuk menghasilkan kerja atau
memungkinkan kerja terjadi dalam suatu sistem rekayasa.

Secara operasional, energon dapat diklasifikasikan sekurang-kurangnya ke
dalam lima jenis: 1. \textbf{Energon-1}: energi fisik; 2.
\textbf{Energon-2}: waktu dan perhatian manusia; 3. \textbf{Energon-3}:
skill dan pengetahuan; 4. \textbf{Energon-4}: token bernilai seperti
uang atau insentif; 5. \textbf{Energon-5}: peran, status, atau
keanggotaan dalam lingkungan kerja dan narasi sosial.

\section*{Definisi Core Engine}\label{definisi-core-engine}
\addcontentsline{toc}{section}{Definisi Core Engine}

\markright{Definisi Core Engine}

\textbf{Core engine} adalah konfigurasi instrumen yang mengubah,
menyalurkan, menyimpan, mengendalikan, atau menstabilkan energi,
material, dan informasi untuk menghasilkan kerja atau keluaran tertentu.

\section*{Definisi Smart Engine}\label{definisi-smart-engine}
\addcontentsline{toc}{section}{Definisi Smart Engine}

\markright{Definisi Smart Engine}

\textbf{Smart engine} adalah konfigurasi instrumen cerdas yang
mengonversi energon menjadi kapasitas kerja, perubahan keadaan, atau
keluaran yang dibutuhkan dalam penyelesaian masalah.

\section*{Definisi Mesin PSKVE}\label{definisi-mesin-pskve}
\addcontentsline{toc}{section}{Definisi Mesin PSKVE}

\markright{Definisi Mesin PSKVE}

\textbf{Mesin PSKVE} adalah smart engine yang mengonversi energon
menjadi lima bentuk utama hasil rekayasa, yaitu \textbf{Product,
Service, Knowledge, Value, dan Environment}.

\section*{Definisi Mesin PUDAL}\label{definisi-mesin-pudal}
\addcontentsline{toc}{section}{Definisi Mesin PUDAL}

\markright{Definisi Mesin PUDAL}

\textbf{Mesin PUDAL} adalah mesin agentik yang menjalankan fungsi
\textbf{Perception, Understanding, Decision, Action, and Learning} untuk
mengendalikan dan mengarahkan operasi mesin PSKVE.

\section*{Definisi Identitas Naratif}\label{definisi-identitas-naratif}
\addcontentsline{toc}{section}{Definisi Identitas Naratif}

\markright{Definisi Identitas Naratif}

\textbf{Identitas naratif} adalah konstruksi kisah hidup yang
diinternalisasi oleh seseorang sebagai dasar bagi pemaknaan diri,
orientasi misi, dan tindakan hidupnya di dunia.

\section*{Definisi DNA Desain}\label{definisi-dna-desain}
\addcontentsline{toc}{section}{Definisi DNA Desain}

\markright{Definisi DNA Desain}

\textbf{DNA desain} adalah representasi sistematis dari pola-pola
desain, skills, aturan, dan struktur konseptual yang memungkinkan suatu
artefak TISE direplikasi, dikembangkan, atau ``dikawinkan'' dengan
artefak lain untuk melahirkan generasi artefak baru.

\section*{Definisi Stakeholder
Bermisi}\label{definisi-stakeholder-bermisi}
\addcontentsline{toc}{section}{Definisi Stakeholder Bermisi}

\markright{Definisi Stakeholder Bermisi}

\textbf{Stakeholder bermisi} adalah pihak manusiawi atau institusional
yang memiliki tujuan, panggilan, atau peran bermakna yang ingin
dijalankan, serta dapat berkontribusi, menerima manfaat, atau mengalami
hambatan dalam suatu konfigurasi solusi.

\section*{Definisi Matrix of Mission}\label{definisi-matrix-of-mission}
\addcontentsline{toc}{section}{Definisi Matrix of Mission}

\markright{Definisi Matrix of Mission}

\textbf{Matrix of Mission} adalah mesin konfigurasi stakeholder yang
memetakan misi, benefit, dan kontribusi timbal balik antar stakeholder
untuk membentuk dan mengevaluasi skema solusi.

\section*{Definisi MOS-7}\label{definisi-mos-7}
\addcontentsline{toc}{section}{Definisi MOS-7}

\markright{Definisi MOS-7}

\textbf{MOS-7} adalah siklus operasional TISE level 5 yang terdiri dari
tujuh langkah: mendengar keluhan, memodelkan stakeholder, mendesain
konfigurasi solusi, membangun komitmen, membangun artefak,
mengoperasionalkan solusi, serta mengevaluasi dan mengaudit integritas
serta dampaknya.

\section*{Definisi TISE level 5}\label{definisi-tise-level-5}
\addcontentsline{toc}{section}{Definisi TISE level 5}

\markright{Definisi TISE level 5}

\textbf{TISE level 5} adalah level rekayasa di mana TISE hadir sebagai
pusat keunggulan yang menggunakan Matrix of Mission dan MOS-7 untuk
memobilisasi potensi manusia yang berlimpah guna memecahkan persoalan
kemanusiaan melalui konfigurasi stakeholder bermisi.

\bookmarksetup{startatroot}

\chapter*{Bagian III. Arsitektur TISE level 5 sebagai Buku
DNA}\label{bagian-iii.-arsitektur-tise-level-5-sebagai-buku-dna}
\addcontentsline{toc}{chapter}{Bagian III. Arsitektur TISE level 5
sebagai Buku DNA}

\markboth{Bagian III. Arsitektur TISE level 5 sebagai Buku DNA}{Bagian
III. Arsitektur TISE level 5 sebagai Buku DNA}

\bookmarksetup{startatroot}

\chapter*{Bab 5. TISE level 5 sebagai Pusat Keunggulan
Rekayasa}\label{bab-5.-tise-level-5-sebagai-pusat-keunggulan-rekayasa}
\addcontentsline{toc}{chapter}{Bab 5. TISE level 5 sebagai Pusat
Keunggulan Rekayasa}

\markboth{Bab 5. TISE level 5 sebagai Pusat Keunggulan Rekayasa}{Bab 5.
TISE level 5 sebagai Pusat Keunggulan Rekayasa}

\section*{Mengapa Pusat Keunggulan?}\label{mengapa-pusat-keunggulan}
\addcontentsline{toc}{section}{Mengapa Pusat Keunggulan?}

\markright{Mengapa Pusat Keunggulan?}

Masalah kemanusiaan bersifat lintas domain, lintas institusi, dan lintas
stakeholder. Karena itu, penyelesaiannya tidak cukup ditangani oleh satu
disiplin, satu unit teknis, atau satu artefak tunggal. Diperlukan
organisasi yang mampu secara sadar mendengar keluhan manusia, memodelkan
konfigurasi stakeholder, merancang solusi, membangun komitmen, dan
menjaga operasi serta integritas solusi secara berkelanjutan.

\section*{Visi}\label{visi}
\addcontentsline{toc}{section}{Visi}

\markright{Visi}

Visi TISE level 5 adalah \textbf{terbentuknya masyarakat yang mampu
memobilisasi potensi manusianya secara cerdas, bermartabat, dan
terorganisasi untuk memecahkan persoalan kemanusiaan dari manusia untuk
manusia}.

\section*{Misi}\label{misi}
\addcontentsline{toc}{section}{Misi}

\markright{Misi}

Misi utama TISE level 5 adalah \textbf{memobilisasi potensi manusia yang
berlimpah untuk memecahkan persoalan kemanusiaan}.

\section*{Prinsip Direktif}\label{prinsip-direktif}
\addcontentsline{toc}{section}{Prinsip Direktif}

\markright{Prinsip Direktif}

Prinsip direktif TISE level 5 adalah \textbf{dari manusia untuk
manusia}.

\section*{Mesin Utama}\label{mesin-utama}
\addcontentsline{toc}{section}{Mesin Utama}

\markright{Mesin Utama}

Mesin utama TISE level 5 adalah \textbf{Matrix of Mission}.

\section*{Unit Analisis Utama}\label{unit-analisis-utama}
\addcontentsline{toc}{section}{Unit Analisis Utama}

\markright{Unit Analisis Utama}

Unit analisis utama TISE level 5 adalah \textbf{stakeholder bermisi
dalam konfigurasi relasional}.

\section*{Artefak Utama}\label{artefak-utama}
\addcontentsline{toc}{section}{Artefak Utama}

\markright{Artefak Utama}

Artefak utama pada level ini adalah \textbf{organisasi, konsorsium,
komunitas stakeholder, atau pusat keunggulan} yang dirancang untuk
menjalankan MOS-7.

\bookmarksetup{startatroot}

\chapter*{Bab 6. Matrix of Mission dan Tipologi
Stakeholder}\label{bab-6.-matrix-of-mission-dan-tipologi-stakeholder}
\addcontentsline{toc}{chapter}{Bab 6. Matrix of Mission dan Tipologi
Stakeholder}

\markboth{Bab 6. Matrix of Mission dan Tipologi Stakeholder}{Bab 6.
Matrix of Mission dan Tipologi Stakeholder}

\section*{Stakeholder sebagai Subjek
Rekayasa}\label{stakeholder-sebagai-subjek-rekayasa}
\addcontentsline{toc}{section}{Stakeholder sebagai Subjek Rekayasa}

\markright{Stakeholder sebagai Subjek Rekayasa}

Dalam TISE level 5, stakeholder adalah aktor utama. Mereka bukan hanya
penerima manfaat, tetapi pembawa misi, pencipta nilai, pengguna artefak,
dan sumber pembelajaran.

\section*{Tipologi Stakeholder}\label{tipologi-stakeholder}
\addcontentsline{toc}{section}{Tipologi Stakeholder}

\markright{Tipologi Stakeholder}

Untuk membuat rekayasa stakeholder menjadi lebih operasional, TISE level
5 mengelompokkan stakeholder ke dalam empat kategori utama, yaitu
\textbf{USER}, \textbf{SOURCE}, \textbf{REGULATOR}, dan
\textbf{PROVIDER}.

\section*{Matrix of Mission sebagai Instrumen Rekayasa
Stakeholder}\label{matrix-of-mission-sebagai-instrumen-rekayasa-stakeholder}
\addcontentsline{toc}{section}{Matrix of Mission sebagai Instrumen
Rekayasa Stakeholder}

\markright{Matrix of Mission sebagai Instrumen Rekayasa Stakeholder}

Secara umum, bentuk matriks dapat digambarkan sebagaimana dirangkum pada
\textbf{Tabel 6.1}.

{\def\LTcaptype{none} % do not increment counter
\begin{longtable}[]{@{}
  >{\raggedright\arraybackslash}p{(\linewidth - 8\tabcolsep) * \real{0.2000}}
  >{\raggedright\arraybackslash}p{(\linewidth - 8\tabcolsep) * \real{0.2000}}
  >{\raggedright\arraybackslash}p{(\linewidth - 8\tabcolsep) * \real{0.2000}}
  >{\raggedright\arraybackslash}p{(\linewidth - 8\tabcolsep) * \real{0.2000}}
  >{\raggedright\arraybackslash}p{(\linewidth - 8\tabcolsep) * \real{0.2000}}@{}}
\toprule\noalign{}
\begin{minipage}[b]{\linewidth}\raggedright
\end{minipage} & \begin{minipage}[b]{\linewidth}\raggedright
A / USER
\end{minipage} & \begin{minipage}[b]{\linewidth}\raggedright
B / SOURCE
\end{minipage} & \begin{minipage}[b]{\linewidth}\raggedright
C / REGULATOR
\end{minipage} & \begin{minipage}[b]{\linewidth}\raggedright
D / PROVIDER
\end{minipage} \\
\midrule\noalign{}
\endhead
\bottomrule\noalign{}
\endlastfoot
1 / USER & Misi USER; Benefit USER & Kontribusi USER ke SOURCE &
Kontribusi USER ke REGULATOR & Kontribusi USER ke PROVIDER \\
2 / SOURCE & Kontribusi SOURCE ke USER & Misi SOURCE; Benefit SOURCE &
Kontribusi SOURCE ke REGULATOR & Kontribusi SOURCE ke PROVIDER \\
3 / REGULATOR & Kontribusi REGULATOR ke USER & Kontribusi REGULATOR ke
SOURCE & Misi REGULATOR; Benefit REGULATOR & Kontribusi REGULATOR ke
PROVIDER \\
4 / PROVIDER & Kontribusi PROVIDER ke USER & Kontribusi PROVIDER ke
SOURCE & Kontribusi PROVIDER ke REGULATOR & Misi PROVIDER; Benefit
PROVIDER \\
\end{longtable}
}

\bookmarksetup{startatroot}

\chapter*{Bab 7. MOS-7 sebagai Siklus
Operasional}\label{bab-7.-mos-7-sebagai-siklus-operasional}
\addcontentsline{toc}{chapter}{Bab 7. MOS-7 sebagai Siklus Operasional}

\markboth{Bab 7. MOS-7 sebagai Siklus Operasional}{Bab 7. MOS-7 sebagai
Siklus Operasional}

\section*{Prinsip Umum}\label{prinsip-umum}
\addcontentsline{toc}{section}{Prinsip Umum}

\markright{Prinsip Umum}

MOS-7 dalam TISE level 5 diposisikan sebagai \textbf{siklus rekayasa
stakeholder bermisi}. Setiap putaran MOS-7 mengubah keluhan kemanusiaan
menjadi konfigurasi stakeholder, konfigurasi stakeholder menjadi artefak
solusi, dan artefak solusi menjadi praktik operasional yang diaudit dan
dipelajari kembali.

\section*{Tujuh Langkah MOS-7}\label{tujuh-langkah-mos-7}
\addcontentsline{toc}{section}{Tujuh Langkah MOS-7}

\markright{Tujuh Langkah MOS-7}

\begin{enumerate}
\def\labelenumi{\arabic{enumi}.}
\tightlist
\item
  \textbf{Identify / Hear}
\item
  \textbf{Model / Discern}
\item
  \textbf{Design / Formulate}
\item
  \textbf{Commit / Seek}
\item
  \textbf{Co-Create / Build}
\item
  \textbf{Operate / Delivery}
\item
  \textbf{Evaluate-Audit / Protect}
\end{enumerate}

\bookmarksetup{startatroot}

\chapter*{Bab 8. Workflow Rekayasa TISE level 5
Full-Scale}\label{bab-8.-workflow-rekayasa-tise-level-5-full-scale}
\addcontentsline{toc}{chapter}{Bab 8. Workflow Rekayasa TISE level 5
Full-Scale}

\markboth{Bab 8. Workflow Rekayasa TISE level 5 Full-Scale}{Bab 8.
Workflow Rekayasa TISE level 5 Full-Scale}

\section*{Fase-Fase Utama Workflow
Rekayasa}\label{fase-fase-utama-workflow-rekayasa}
\addcontentsline{toc}{section}{Fase-Fase Utama Workflow Rekayasa}

\markright{Fase-Fase Utama Workflow Rekayasa}

Secara linear, workflow rekayasa TISE level 5 dapat dibagi ke dalam enam
fase utama: 1. \textbf{Persepsi Masalah} 2. \textbf{Konsep dan Model} 3.
\textbf{Desain dan Planning} 4. \textbf{Konstruksi} 5.
\textbf{Operasional} 6. \textbf{Evaluasi}

Untuk memperjelas hubungan antara fase linear, putaran MOS-7, dan
penurunan level rekayasa dari TISE level 5 ke TISE level 0,
\textbf{Tabel 8.1} merangkum workflow rekayasa TISE level 5 secara
full-scale.

{\def\LTcaptype{none} % do not increment counter
\begin{longtable}[]{@{}
  >{\raggedright\arraybackslash}p{(\linewidth - 8\tabcolsep) * \real{0.2000}}
  >{\raggedright\arraybackslash}p{(\linewidth - 8\tabcolsep) * \real{0.2000}}
  >{\raggedright\arraybackslash}p{(\linewidth - 8\tabcolsep) * \real{0.2000}}
  >{\raggedright\arraybackslash}p{(\linewidth - 8\tabcolsep) * \real{0.2000}}
  >{\raggedright\arraybackslash}p{(\linewidth - 8\tabcolsep) * \real{0.2000}}@{}}
\toprule\noalign{}
\begin{minipage}[b]{\linewidth}\raggedright
Fase Linear
\end{minipage} & \begin{minipage}[b]{\linewidth}\raggedright
Fokus Utama
\end{minipage} & \begin{minipage}[b]{\linewidth}\raggedright
Putaran MOS-7 yang Dominan
\end{minipage} & \begin{minipage}[b]{\linewidth}\raggedright
Level TISE yang Paling Dominan
\end{minipage} & \begin{minipage}[b]{\linewidth}\raggedright
Keluaran Utama
\end{minipage} \\
\midrule\noalign{}
\endhead
\bottomrule\noalign{}
\endlastfoot
Persepsi Masalah & Mendengar keluhan manusia, mengidentifikasi
stakeholder utama, peluang provider, dan alternatif konfigurasi & Hear,
Model & TISE level 5 & Rumusan persoalan kemanusiaan, daftar
stakeholder, alternatif Matrix of Mission \\
Konsep dan Model & Membandingkan konfigurasi stakeholder, memilih
konfigurasi terbaik, merumuskan arsitektur solusi & Hear, Model, Design,
Commit & TISE level 5 dan TISE level 4 & Matrix of Mission terpilih,
konsep platform, arah artefak anak \\
Desain dan Planning & Menurunkan kebutuhan dari level stakeholder ke
level DNA desain, naratif, agentik, PSKVE, dan core engineering &
Design, Commit, Build & TISE level 4, TISE level 3, TISE level 2, TISE
level 1 & DNA platform, DNA gagasan baru, desain teater naratif,
rancangan PUDAL, rancangan PSKVE, spesifikasi elemen TISE level 0 \\
Konstruksi & Membangun atau memodifikasi artefak anak, platform, agen,
layanan, dan elemen teknis yang diperlukan & Build, Commit, Operate &
TISE level 4 sampai TISE level 0 & Artefak anak, modul layanan,
mekanisme operasi, elemen teknis siap pakai \\
Operasional & Menjalankan solusi dalam kehidupan nyata dengan peran
USER, SOURCE, REGULATOR, dan PROVIDER & Operate, Evaluate & Seluruh
level bekerja bersama, dipandu TISE level 5 & Solusi hidup, operasi
provider, pengalaman user, legitimasi regulator, value exchange nyata \\
Evaluasi & Audit integritas, dampak, keberlanjutan, fairness, dan
peluang redesign & Evaluate, Hear, Model & TISE level 5 & Hasil
evaluasi, audit, pembelajaran, redesign, dan persiapan siklus
berikutnya \\
\end{longtable}
}

\section*{Workflow Hipotetikal
Full-Scale}\label{workflow-hipotetikal-full-scale}
\addcontentsline{toc}{section}{Workflow Hipotetikal Full-Scale}

\markright{Workflow Hipotetikal Full-Scale}

\begin{enumerate}
\def\labelenumi{\arabic{enumi}.}
\tightlist
\item
  Menyiapkan TISE level 5 dan MOS-7.
\item
  Mengembangkan TISE level 4 yang relevan.
\item
  Mengembangkan TISE level 3 dalam perspektif TISE level 4.
\item
  Mengembangkan TISE level 2 dalam konteks TISE level 3.
\item
  Mengembangkan TISE level 1 dalam konteks TISE level 2.
\item
  Mengembangkan TISE level 0.
\end{enumerate}

\bookmarksetup{startatroot}

\chapter*{Bagian IV. Domain Aplikasi dan Reusability
DNA}\label{bagian-iv.-domain-aplikasi-dan-reusability-dna}
\addcontentsline{toc}{chapter}{Bagian IV. Domain Aplikasi dan
Reusability DNA}

\markboth{Bagian IV. Domain Aplikasi dan Reusability DNA}{Bagian IV.
Domain Aplikasi dan Reusability DNA}

\bookmarksetup{startatroot}

\chapter*{Bab 9. Persoalan Kemanusiaan sebagai Objek
Rekayasa}\label{bab-9.-persoalan-kemanusiaan-sebagai-objek-rekayasa}
\addcontentsline{toc}{chapter}{Bab 9. Persoalan Kemanusiaan sebagai
Objek Rekayasa}

\markboth{Bab 9. Persoalan Kemanusiaan sebagai Objek Rekayasa}{Bab 9.
Persoalan Kemanusiaan sebagai Objek Rekayasa}

Masalah kemanusiaan adalah situasi di mana manusia atau komunitas
manusia mengalami hambatan dalam menjalankan misinya, kehilangan agensi,
mengalami penderitaan, atau hidup dalam kondisi yang tidak selaras
dengan nilai dan martabat kemanusiaannya.

\bookmarksetup{startatroot}

\chapter*{Bab 10. Domain Aplikasi TISE level
5}\label{bab-10.-domain-aplikasi-tise-level-5}
\addcontentsline{toc}{chapter}{Bab 10. Domain Aplikasi TISE level 5}

\markboth{Bab 10. Domain Aplikasi TISE level 5}{Bab 10. Domain Aplikasi
TISE level 5}

Domain aplikasi TISE level 5 dapat mencakup antara lain: 1. makanan
sehat; 2. lansia berkualitas; 3. pendidikan teknik abad ke-21; 4.
konsentrasi ADHD; 5. decision making masyarakat cerdas; 6. rekayasa
finansial untuk proyek IT; 7. personal finance; 8. digital loan.

\bookmarksetup{startatroot}

\chapter*{Bab 11. Reusability, Modifikasi, dan Rekayasa Berbasis
Template
DNA}\label{bab-11.-reusability-modifikasi-dan-rekayasa-berbasis-template-dna}
\addcontentsline{toc}{chapter}{Bab 11. Reusability, Modifikasi, dan
Rekayasa Berbasis Template DNA}

\markboth{Bab 11. Reusability, Modifikasi, dan Rekayasa Berbasis
Template DNA}{Bab 11. Reusability, Modifikasi, dan Rekayasa Berbasis
Template DNA}

\section*{Buku DNA sebagai Template
Rekayasa}\label{buku-dna-sebagai-template-rekayasa}
\addcontentsline{toc}{section}{Buku DNA sebagai Template Rekayasa}

\markright{Buku DNA sebagai Template Rekayasa}

Pada akhirnya, monograf ini dimaksudkan sebagai \textbf{Buku DNA} bagi
TISE level 5. Artinya, buku ini tidak hanya menjelaskan teori, tetapi
menyediakan pola dasar yang dapat digunakan ulang oleh rekayasawan,
stakeholder, dan agents dalam domain yang berbeda-beda.

\section*{Elemen Stabil dalam Buku
DNA}\label{elemen-stabil-dalam-buku-dna}
\addcontentsline{toc}{section}{Elemen Stabil dalam Buku DNA}

\markright{Elemen Stabil dalam Buku DNA}

Elemen-elemen yang relatif stabil dan sebaiknya dipertahankan sebagai
template umum meliputi: 1. hirarki level rekayasa dari TISE level 0
hingga TISE level 5; 2. konsep energon, core engine, smart engine,
PSKVE, PUDAL, identitas naratif, DNA desain, dan Matrix of Mission; 3.
tipologi stakeholder USER, SOURCE, REGULATOR, dan PROVIDER; 4. struktur
MOS-7; 5. workflow full-scale dari level tertinggi ke level terbawah; 6.
prinsip direktif dari manusia untuk manusia.

\section*{Elemen Adaptif yang Dapat
Dimodifikasi}\label{elemen-adaptif-yang-dapat-dimodifikasi}
\addcontentsline{toc}{section}{Elemen Adaptif yang Dapat Dimodifikasi}

\markright{Elemen Adaptif yang Dapat Dimodifikasi}

Bagian-bagian yang dimaksudkan untuk dimodifikasi oleh rekayasawan
sesuai konteks antara lain: 1. domain persoalan kemanusiaan yang
dipilih; 2. identitas stakeholder utama dan stakeholder pendukung; 3.
isi Matrix of Mission; 4. DNA platform yang sudah tersedia dan DNA
gagasan baru yang ditambahkan; 5. artefak anak yang ingin dilahirkan; 6.
model operasi PROVIDER; 7. indikator evaluasi, audit, dan keberhasilan.

\bookmarksetup{startatroot}

\chapter*{Bagian V. Pengembangan Ilmiah dan
Kelembagaan}\label{bagian-v.-pengembangan-ilmiah-dan-kelembagaan}
\addcontentsline{toc}{chapter}{Bagian V. Pengembangan Ilmiah dan
Kelembagaan}

\markboth{Bagian V. Pengembangan Ilmiah dan Kelembagaan}{Bagian V.
Pengembangan Ilmiah dan Kelembagaan}

\bookmarksetup{startatroot}

\chapter*{Bab 12. Proposisi, Agenda Riset, dan Pengembangan
Lanjutan}\label{bab-12.-proposisi-agenda-riset-dan-pengembangan-lanjutan}
\addcontentsline{toc}{chapter}{Bab 12. Proposisi, Agenda Riset, dan
Pengembangan Lanjutan}

\markboth{Bab 12. Proposisi, Agenda Riset, dan Pengembangan
Lanjutan}{Bab 12. Proposisi, Agenda Riset, dan Pengembangan Lanjutan}

\section*{Proposisi Inti TISE level
5}\label{proposisi-inti-tise-level-5}
\addcontentsline{toc}{section}{Proposisi Inti TISE level 5}

\markright{Proposisi Inti TISE level 5}

\begin{enumerate}
\def\labelenumi{\arabic{enumi}.}
\tightlist
\item
  \textbf{Persoalan kemanusiaan abad ke-21 pada umumnya adalah persoalan
  konfigurasi stakeholder, bukan semata persoalan artefak tunggal.}
\item
  \textbf{Karena itu, mesin utama pada level rekayasa tertinggi bukan
  hanya mesin teknis, melainkan mesin konfigurasi misi, yaitu Matrix of
  Mission.}
\item
  \textbf{Keberhasilan solusi tidak hanya ditentukan oleh performa
  teknis artefak, tetapi oleh kesesuaian misi, benefit, kontribusi, dan
  komitmen antar stakeholder.}
\item
  \textbf{MOS-7 diperlukan sebagai siklus operasional untuk membentuk,
  menggerakkan, dan mengevaluasi konfigurasi stakeholder tersebut.}
\item
  \textbf{TISE level 5 diperlukan sebagai bentuk kelembagaan yang
  memungkinkan mobilisasi potensi manusia secara sistematis untuk
  memecahkan persoalan kemanusiaan.}
\end{enumerate}

\bookmarksetup{startatroot}

\chapter*{Bab 13. Kesimpulan: Buku DNA untuk Rekayasawan, Stakeholder,
dan
Agents}\label{bab-13.-kesimpulan-buku-dna-untuk-rekayasawan-stakeholder-dan-agents}
\addcontentsline{toc}{chapter}{Bab 13. Kesimpulan: Buku DNA untuk
Rekayasawan, Stakeholder, dan Agents}

\markboth{Bab 13. Kesimpulan: Buku DNA untuk Rekayasawan, Stakeholder,
dan Agents}{Bab 13. Kesimpulan: Buku DNA untuk Rekayasawan, Stakeholder,
dan Agents}

TISE level 5 adalah tahap perkembangan paradigma rekayasa cerdas yang
menempatkan TISE sebagai pusat keunggulan untuk memecahkan persoalan
kemanusiaan. Dengan menjadikan MOS-7 sebagai siklus operasional, TISE
level 5 mampu mengubah keluhan manusia menjadi agenda rekayasa, mengubah
kombinasi stakeholder menjadi kekuatan kolektif, dan mengubah artefak
menjadi solusi yang hidup dalam operasi nyata.

Sebagai Buku DNA, monograf ini menyimpan struktur, prinsip, bahasa,
ontologi, workflow, mesin, dan pola-pola desain yang dapat dibaca dan
diacu bersama oleh para rekayasawan, stakeholder, dan agents. Setiap
pihak tidak perlu membaca seluruh sistem dengan intensitas yang sama;
masing-masing dapat membaca bagian yang relevan dengan misi dan
perannya. Tetapi justru karena seluruh pihak mengacu pada buku yang
sama, konsistensi, kesinambungan, dan reproduktibilitas TISE level 5
dapat dijaga.

\bookmarksetup{startatroot}

\chapter*{Daftar Pustaka}\label{daftar-pustaka}
\addcontentsline{toc}{chapter}{Daftar Pustaka}

\markboth{Daftar Pustaka}{Daftar Pustaka}

\protect\phantomsection\label{refs}
\begin{CSLReferences}{1}{1}
\bibitem[\citeproctext]{ref-berkovich2014}
Berkovich, Marina, Jan Marco Leimeister, Axel Hoffmann, and Helmut
Krcmar. 2014. {``A Requirements Data Model for Product Service
Systems.''} \emph{Requirements Engineering} 19: 161--86.
\url{https://doi.org/10.1007/s00766-012-0164-1}.

\bibitem[\citeproctext]{ref-giordano2018}
Giordano, Francesca, Nicola Morelli, Amalia De Götzen, and Jonas
Hunziker. 2018. {``The Stakeholder Map: A Conversation Tool for
Designing People-Led Public Services.''} \emph{Proceedings of the
ServDes.2018 Conference}, Linköping electronic conference proceedings,
vol. 150: 582--97.

\bibitem[\citeproctext]{ref-lievens2021}
Lievens, Annouk, and Vera Blazevic. 2021. {``A Service Design
Perspective on the Stakeholder Engagement Journey During B2B Innovation:
Challenges and Future Research Agenda.''} \emph{Industrial Marketing
Management} 95: 128--41.
\url{https://doi.org/10.1016/j.indmarman.2021.04.007}.

\bibitem[\citeproctext]{ref-morelli2002}
Morelli, Nicola. 2002. {``Designing Product/Service Systems: A
Methodological Exploration.''} \emph{Design Issues} 18 (3): 3--17.
\url{https://doi.org/10.1162/074793602320223253}.

\bibitem[\citeproctext]{ref-sholihah2019}
Sholihah, M., Y. Mitake, T. Nakada, and Y. Shimomura. 2019.
{``Innovative Design Method for a Valuable Product-Service System:
Concretizing Multi-Stakeholder Requirements.''} \emph{Journal of
Advanced Mechanical Design, Systems and Manufacturing} 13 (5): 1--12.
\url{https://doi.org/10.1299/jamdsm.2019jamdsm0091}.

\bibitem[\citeproctext]{ref-vezzoli2014}
Vezzoli, Carlo, Cindy Kohtala, Amrit Srinivasan, et al. 2014.
\emph{Product-Service System Design for Sustainability}. Greenleaf
Publishing. \url{https://doi.org/10.4324/9781351278003}.

\end{CSLReferences}

\bookmarksetup{startatroot}

\chapter{Introduction}\label{introduction}

This is a book created from markdown and executable code.

See (\textbf{knuth84?}) for additional discussion of literate
programming.

\bookmarksetup{startatroot}

\chapter{Summary}\label{summary}

In summary, this book has no content whatsoever.

\bookmarksetup{startatroot}

\chapter*{References}\label{references}
\addcontentsline{toc}{chapter}{References}

\markboth{References}{References}

\protect\phantomsection\label{refs}
\begin{CSLReferences}{1}{1}
\bibitem[\citeproctext]{ref-berkovich2014}
Berkovich, Marina, Jan Marco Leimeister, Axel Hoffmann, and Helmut
Krcmar. 2014. {``A Requirements Data Model for Product Service
Systems.''} \emph{Requirements Engineering} 19: 161--86.
\url{https://doi.org/10.1007/s00766-012-0164-1}.

\bibitem[\citeproctext]{ref-giordano2018}
Giordano, Francesca, Nicola Morelli, Amalia De Götzen, and Jonas
Hunziker. 2018. {``The Stakeholder Map: A Conversation Tool for
Designing People-Led Public Services.''} \emph{Proceedings of the
ServDes.2018 Conference}, Linköping electronic conference proceedings,
vol. 150: 582--97.

\bibitem[\citeproctext]{ref-lievens2021}
Lievens, Annouk, and Vera Blazevic. 2021. {``A Service Design
Perspective on the Stakeholder Engagement Journey During B2B Innovation:
Challenges and Future Research Agenda.''} \emph{Industrial Marketing
Management} 95: 128--41.
\url{https://doi.org/10.1016/j.indmarman.2021.04.007}.

\bibitem[\citeproctext]{ref-morelli2002}
Morelli, Nicola. 2002. {``Designing Product/Service Systems: A
Methodological Exploration.''} \emph{Design Issues} 18 (3): 3--17.
\url{https://doi.org/10.1162/074793602320223253}.

\bibitem[\citeproctext]{ref-sholihah2019}
Sholihah, M., Y. Mitake, T. Nakada, and Y. Shimomura. 2019.
{``Innovative Design Method for a Valuable Product-Service System:
Concretizing Multi-Stakeholder Requirements.''} \emph{Journal of
Advanced Mechanical Design, Systems and Manufacturing} 13 (5): 1--12.
\url{https://doi.org/10.1299/jamdsm.2019jamdsm0091}.

\bibitem[\citeproctext]{ref-vezzoli2014}
Vezzoli, Carlo, Cindy Kohtala, Amrit Srinivasan, et al. 2014.
\emph{Product-Service System Design for Sustainability}. Greenleaf
Publishing. \url{https://doi.org/10.4324/9781351278003}.

\end{CSLReferences}




\end{document}
